\documentclass[a4paper]{article}

\usepackage{graphicx}
\usepackage{amsmath,amssymb}
\usepackage{minted}
\usepackage{multirow}
\usepackage{float}
\usepackage{algorithm}
\usepackage{algpseudocode}
\usepackage{todonotes}
\usepackage{forest}
\usepackage{xstring}
\usepackage{xcolor}
\usepackage[most]{tcolorbox}
\usepackage{xparse}
\usepackage{natbib}

\usetikzlibrary{graphs,graphdrawing}
\usegdlibrary{layered}

\title{Strict recursive Preview workspace path resolution}
\date{2026-09-07}
\author{Mohammad Rahmani}

\begin{document}
\maketitle

\section{Problem}
A stale running Mind Mirror process could still use the older recursive input
resolver. That resolver searched the repository for suffix matches when a
relative \verb|\input| or \verb|\include| was not found beside its referring
file. A thesis Preview could therefore pull an unrelated file from
\verb|src/kind/presentation| even though the thesis never referenced that
presentation tree.

\section{Canonical behavior}
The canonical workspace recursive resolver does not perform repository-wide
suffix search. Relative paths are resolved only beside the referring file.
Paths beginning with \verb|/| are resolved below
\verb|search.default_path| first and then as true operating-system absolute
paths.

The application launches \verb|EnhancedMainWindow|, which configures Preview
with \verb|WorkspaceRecursiveInputResolver| and refreshes the current Preview.

\section{Regression coverage}
Regression tests now verify that:
\begin{itemize}
   \item a missing thesis-relative input never selects an unrelated
   presentation file with a matching suffix;
   \item a leading-slash workspace reference resolves below
   \verb|search.default_path|;
   \item the application launch path wires the workspace recursive resolver.
\end{itemize}

\end{document}
